% Generated draft: requires mathematical and source review.
\documentclass[11pt]{article}
\usepackage[margin=25mm]{geometry}
\usepackage{fontspec}
\setmainfont{Latin Modern Roman}
\usepackage{amsmath,amssymb,mathtools,unicode-math}
\setmathfont{Latin Modern Math}
\usepackage{enumitem,needspace,xurl,hyperref}
\setlength{\parindent}{0pt}
\setlength{\parskip}{5pt}
\setlength{\emergencystretch}{3em}
\newcommand{\N}{\mathbb{N}}
\newcommand{\Z}{\mathbb{Z}}
\newcommand{\Q}{\mathbb{Q}}
\newcommand{\R}{\mathbb{R}}
\newcommand{\C}{\mathbb{C}}
\newcommand{\F}{\mathbb{F}}
\newcommand{\A}{\mathbb{A}}
\newcommand{\PP}{\mathbb{P}}
\newcommand{\E}{\mathbb{E}}
\newcommand{\eps}{\varepsilon}
\newcommand{\dd}{\,\mathrm{d}}
\begin{document}
\section*{UnsolvedMath ID: 700548}
\textbf{Source problem code:} catalogue entry
\subsection{Definitions and mathematical statement}
A group \(G\) is soluble if its derived series, defined by \(G^{(0)}=G\) and \(G^{(i+1)}=[G^{(i)},G^{(i)}]\), satisfies \(G^{(d)}=1\) for some integer \(d\geq 0\). It is torsion-free if \(g^n=1\), for \(g\in G\) and an integer \(n\geq 1\), implies \(g=1\).

For this draft, the weak minimum condition for normal subgroups means that every descending chain
\[
N_0\geq N_1\geq N_2\geq\cdots,\qquad N_i\trianglelefteq G,
\]
has an index \(i_0\) such that \([N_i:N_{i+1}]<\infty\) for every \(i\geq i_0\). Equivalently, there is no infinite descending chain of normal subgroups of \(G\) in which every successive index is infinite. Normality here is in the ambient group \(G\).

A soluble group \(G\) is minimax if it has a finite subnormal series
\[
1=G_0\trianglelefteq G_1\trianglelefteq\cdots\trianglelefteq G_r=G
\]
whose factors are finite cyclic groups, infinite cyclic groups, or quasicyclic groups \(C_{p^\infty}\), where the prime \(p\) may vary between factors. Here \(C_{p^\infty}\) is the direct limit of cyclic groups of orders \(p,p^2,p^3,\ldots\) under the usual inclusions. Torsion-freeness is imposed on \(G\), and does not exclude torsion in its sections.

\textbf{Question (D. I. Zaitsev, 1984).} Is every soluble torsion-free group satisfying the weak minimum condition for normal subgroups minimax?

The supplied record does not define either the weak minimum condition or minimax. The explicit conventions above give the intended statement a precise reading, but their agreement with the cited source requires human verification.
\subsection{Short English statement}
Must every soluble torsion-free group whose descending chains of normal subgroups eventually have only finite successive indices admit a finite subnormal series with finite cyclic, infinite cyclic, or quasicyclic factors? This question is attributed to D. I. Zaitsev (1984). The supplied registry labels it open; that label has not been independently verified.
\subsection{Sources}
\begin{itemize}
\item Local UnsolvedMath registry, supplied complete source record: permanent ID \(700548\), exact source code \texttt{catalogue entry}, title \emph{Kourovka Notebook Problem 9.13}. The record supplies the status \texttt{open}.
\item \emph{The Kourovka Notebook (Unsolved Problems in Group Theory)}, 9th issue (1984), Problem 9.13; proposed by D. I. Zaitsev, according to the supplied attribution.
\item The supplied background also identifies the 21st issue edition (2026), \texttt{arXiv:1401.0300v48}. The supplied collection description names E. I. Khukhro and V. D. Mazurov as editors. These bibliographic details are supplied source information and have not been independently checked.
\end{itemize}
\end{document}
